\documentclass[a4paper,11pt,fontset=none]{ctexart}
% ============================================================
%  hjstyle.tex  —  shared preamble for the IMP/HIAF application
%  Use:  \documentclass[a4paper,11pt,fontset=none]{ctexart}
%        % ============================================================
%  hjstyle.tex  —  shared preamble for the IMP/HIAF application
%  Use:  \documentclass[a4paper,11pt,fontset=none]{ctexart}
%        % ============================================================
%  hjstyle.tex  —  shared preamble for the IMP/HIAF application
%  Use:  \documentclass[a4paper,11pt,fontset=none]{ctexart}
%        \input{hjstyle.tex}
%  Compile with: xelatex
% ============================================================

\usepackage{geometry}
\geometry{a4paper, top=2.4cm, bottom=2.3cm, left=2.4cm, right=2.4cm, headsep=10pt}

\usepackage{amsmath,amssymb}
\usepackage{bm}
\usepackage{graphicx}
\usepackage{wrapfig}
\usepackage{float}
\usepackage{array}
\usepackage{booktabs}
\usepackage{tabularx}
\usepackage{enumitem}
\usepackage[dvipsnames,table]{xcolor}
\usepackage{titlesec}
\usepackage{fancyhdr}
\usepackage[normalem]{ulem}
\usepackage{caption}
\usepackage{tikz}
\usetikzlibrary{arrows.meta,positioning,calc,shapes.geometric,backgrounds,fit}
\usepackage{fontspec}

% ---------- Fonts: portable across machines (see README.md) ----------
% Latin = XCharter (bundled with TeX Live). CJK = macOS system fonts when present
% (preserves the original look), else Fandol (bundled with TeX Live) so the document
% builds on ANY TeX Live install (Mac, Linux, Windows) with no extra font setup.
\setmainfont{XCharter}% Latin: XCharter ships with TeX Live (always available)
\IfFontExistsTF{Songti SC}
  {\setCJKmainfont{Songti SC}\newCJKfontfamily\song{Songti SC}}
  {\setCJKmainfont{FandolSong-Regular.otf}[BoldFont=FandolSong-Bold.otf]%
   \newCJKfontfamily\song{FandolSong-Regular.otf}}
\IfFontExistsTF{Heiti SC}
  {\setCJKsansfont{Heiti SC}\newCJKfontfamily\hei{Heiti SC}}% CJK sans for emphasis/headings
  {\setCJKsansfont{FandolHei-Regular.otf}[BoldFont=FandolHei-Bold.otf]%
   \newCJKfontfamily\hei{FandolHei-Regular.otf}}
% route enclosed numbers ①-⑳ and Roman numerals to the CJK font (XCharter lacks them)
\xeCJKDeclareCharClass{CJK}{"2160 -> "2188}
\xeCJKDeclareCharClass{CJK}{"2460 -> "24FF}

% ---------- Colour palette ----------
\definecolor{impblue}{RGB}{31,73,125}     % deep institutional blue
\definecolor{impsky}{RGB}{46,116,181}     % brighter accent blue
\definecolor{impgray}{RGB}{120,120,120}   % header/footer grey
\definecolor{impink}{RGB}{20,40,70}        % near-black ink for key text
\definecolor{impbg}{RGB}{238,243,249}      % pale blue panel background

% ---------- Paragraph style: clean block layout ----------
\setlength{\parindent}{0pt}
\setlength{\parskip}{0.55em}
\linespread{1.06}
\emergencystretch=1.6em   % absorb long unbreakable Latin/math runs, avoid overfull
\setlist{nosep, leftmargin=1.5em, topsep=2pt, itemsep=3pt}

% ---------- Section headings (blue, ruled) ----------
\titleformat{\section}[hang]
  {\Large\bfseries\color{impblue}}{\thesection.}{0.5em}{}
  [{\vspace{1pt}\color{impblue!35}\titlerule[1pt]}]
\titlespacing*{\section}{0pt}{1.1em}{0.5em}

\titleformat{\subsection}[hang]
  {\large\bfseries\color{impsky}}{\thesubsection}{0.5em}{}
\titlespacing*{\subsection}{0pt}{0.8em}{0.25em}

\titleformat{\subsubsection}[runin]
  {\bfseries\color{impink}}{}{0em}{}[\,]
\titlespacing*{\subsubsection}{0pt}{0.4em}{0.5em}

% ---------- Captions ----------
\captionsetup{font=small, labelfont={bf,color=impblue}, labelsep=period, skip=4pt}

% ---------- Running header / footer ----------
\newcommand{\doctitle}{}
\pagestyle{fancy}
\fancyhf{}
\setlength{\headheight}{14pt}
\fancyhead[L]{\small\color{impgray}\,马越\ \textperiodcentered\ Yue Ma}
\fancyhead[R]{\small\color{impgray}\doctitle\,}
\fancyfoot[C]{\small\color{impgray}\thepage}
\renewcommand{\headrulewidth}{0.4pt}
\renewcommand{\footrulewidth}{0pt}
\renewcommand{\headrule}{\hbox to\headwidth{\color{impblue!30}\leaders\hrule height \headrulewidth\hfill}}

% ---------- Physics macros ----------
\newcommand{\hn}[2]{\ensuremath{^{#1}_{\Lambda}\mathrm{#2}}}        % single-Lambda hypernucleus  \hn{3}{H}
\newcommand{\hnn}[2]{\ensuremath{^{#1}_{\Lambda\Lambda}\mathrm{#2}}}% double-Lambda hypernucleus  \hnn{6}{He}
\newcommand{\hsig}[2]{\ensuremath{^{#1}_{\Sigma}\mathrm{#2}}}       % Sigma hypernucleus  \hsig{3}{H}
\newcommand{\dbar}{\ensuremath{\bar{d}}}
\newcommand{\pbar}{\ensuremath{\bar{p}}}
\newcommand{\Kbar}{\ensuremath{\bar{K}}}
\newcommand{\KN}{\ensuremath{\bar{K}N}}
\newcommand{\Kmpp}{\ensuremath{K^{-}pp}}
\newcommand{\KNN}{\ensuremath{\bar{K}NN}}
\newcommand{\Lam}{\ensuremath{\Lambda}}
\newcommand{\Sig}{\ensuremath{\Sigma}}
\newcommand{\Xibar}{\ensuremath{\Xi}}
\newcommand{\Lc}{\ensuremath{\Lambda_c^{+}}}
\newcommand{\gevu}{\,GeV/u}
\newcommand{\mevc}{\,MeV/$c$}

% ---------- Emphasis helpers ----------
\newcommand{\lead}[1]{\textcolor{impblue}{\textbf{#1}}}   % blue bold lead-in
\newcommand{\keypt}[1]{\textbf{#1}}                         % key statement
\newcommand{\keyu}[1]{\uline{\textbf{#1}}}                  % strongest emphasis (use sparingly)
\newcommand{\role}[1]{\textcolor{impsky}{\textbf{[#1]}}}    % authorship/role tag
\newcommand{\tech}[1]{\textmd{#1}}                          % technical term/acronym: keep normal weight even inside bold

% ---------- Blue-bar callout block (for "my contribution" etc.) ----------
\newenvironment{callout}
  {\par\smallskip\noindent\begin{list}{}{%
     \setlength{\leftmargin}{1.1em}\setlength{\rightmargin}{0pt}%
     \setlength{\topsep}{2pt}\setlength{\itemsep}{0pt}}%
     \item[]\color{impink}%
     \def\FrameRule{0pt}}
  {\end{list}\par\smallskip}

% Title banner used at the top of each document
\newcommand{\hjbanner}[3]{% #1 = CN title, #2 = EN subtitle, #3 = meta line
  {\noindent\Huge\bfseries\color{impblue} #1\par}
  \vspace{2pt}
  {\noindent\large\color{impink} #2\par}
  \vspace{4pt}
  {\noindent\small\color{impgray} #3\par}
  \vspace{2pt}
  {\noindent\color{impblue!45}\rule{\linewidth}{1.2pt}\par}
  \vspace{4pt}
}

\usepackage[hidelinks]{hyperref}

%  Compile with: xelatex
% ============================================================

\usepackage{geometry}
\geometry{a4paper, top=2.4cm, bottom=2.3cm, left=2.4cm, right=2.4cm, headsep=10pt}

\usepackage{amsmath,amssymb}
\usepackage{bm}
\usepackage{graphicx}
\usepackage{wrapfig}
\usepackage{float}
\usepackage{array}
\usepackage{booktabs}
\usepackage{tabularx}
\usepackage{enumitem}
\usepackage[dvipsnames,table]{xcolor}
\usepackage{titlesec}
\usepackage{fancyhdr}
\usepackage[normalem]{ulem}
\usepackage{caption}
\usepackage{tikz}
\usetikzlibrary{arrows.meta,positioning,calc,shapes.geometric,backgrounds,fit}
\usepackage{fontspec}

% ---------- Fonts: portable across machines (see README.md) ----------
% Latin = XCharter (bundled with TeX Live). CJK = macOS system fonts when present
% (preserves the original look), else Fandol (bundled with TeX Live) so the document
% builds on ANY TeX Live install (Mac, Linux, Windows) with no extra font setup.
\setmainfont{XCharter}% Latin: XCharter ships with TeX Live (always available)
\IfFontExistsTF{Songti SC}
  {\setCJKmainfont{Songti SC}\newCJKfontfamily\song{Songti SC}}
  {\setCJKmainfont{FandolSong-Regular.otf}[BoldFont=FandolSong-Bold.otf]%
   \newCJKfontfamily\song{FandolSong-Regular.otf}}
\IfFontExistsTF{Heiti SC}
  {\setCJKsansfont{Heiti SC}\newCJKfontfamily\hei{Heiti SC}}% CJK sans for emphasis/headings
  {\setCJKsansfont{FandolHei-Regular.otf}[BoldFont=FandolHei-Bold.otf]%
   \newCJKfontfamily\hei{FandolHei-Regular.otf}}
% route enclosed numbers ①-⑳ and Roman numerals to the CJK font (XCharter lacks them)
\xeCJKDeclareCharClass{CJK}{"2160 -> "2188}
\xeCJKDeclareCharClass{CJK}{"2460 -> "24FF}

% ---------- Colour palette ----------
\definecolor{impblue}{RGB}{31,73,125}     % deep institutional blue
\definecolor{impsky}{RGB}{46,116,181}     % brighter accent blue
\definecolor{impgray}{RGB}{120,120,120}   % header/footer grey
\definecolor{impink}{RGB}{20,40,70}        % near-black ink for key text
\definecolor{impbg}{RGB}{238,243,249}      % pale blue panel background

% ---------- Paragraph style: clean block layout ----------
\setlength{\parindent}{0pt}
\setlength{\parskip}{0.55em}
\linespread{1.06}
\emergencystretch=1.6em   % absorb long unbreakable Latin/math runs, avoid overfull
\setlist{nosep, leftmargin=1.5em, topsep=2pt, itemsep=3pt}

% ---------- Section headings (blue, ruled) ----------
\titleformat{\section}[hang]
  {\Large\bfseries\color{impblue}}{\thesection.}{0.5em}{}
  [{\vspace{1pt}\color{impblue!35}\titlerule[1pt]}]
\titlespacing*{\section}{0pt}{1.1em}{0.5em}

\titleformat{\subsection}[hang]
  {\large\bfseries\color{impsky}}{\thesubsection}{0.5em}{}
\titlespacing*{\subsection}{0pt}{0.8em}{0.25em}

\titleformat{\subsubsection}[runin]
  {\bfseries\color{impink}}{}{0em}{}[\,]
\titlespacing*{\subsubsection}{0pt}{0.4em}{0.5em}

% ---------- Captions ----------
\captionsetup{font=small, labelfont={bf,color=impblue}, labelsep=period, skip=4pt}

% ---------- Running header / footer ----------
\newcommand{\doctitle}{}
\pagestyle{fancy}
\fancyhf{}
\setlength{\headheight}{14pt}
\fancyhead[L]{\small\color{impgray}\,马越\ \textperiodcentered\ Yue Ma}
\fancyhead[R]{\small\color{impgray}\doctitle\,}
\fancyfoot[C]{\small\color{impgray}\thepage}
\renewcommand{\headrulewidth}{0.4pt}
\renewcommand{\footrulewidth}{0pt}
\renewcommand{\headrule}{\hbox to\headwidth{\color{impblue!30}\leaders\hrule height \headrulewidth\hfill}}

% ---------- Physics macros ----------
\newcommand{\hn}[2]{\ensuremath{^{#1}_{\Lambda}\mathrm{#2}}}        % single-Lambda hypernucleus  \hn{3}{H}
\newcommand{\hnn}[2]{\ensuremath{^{#1}_{\Lambda\Lambda}\mathrm{#2}}}% double-Lambda hypernucleus  \hnn{6}{He}
\newcommand{\hsig}[2]{\ensuremath{^{#1}_{\Sigma}\mathrm{#2}}}       % Sigma hypernucleus  \hsig{3}{H}
\newcommand{\dbar}{\ensuremath{\bar{d}}}
\newcommand{\pbar}{\ensuremath{\bar{p}}}
\newcommand{\Kbar}{\ensuremath{\bar{K}}}
\newcommand{\KN}{\ensuremath{\bar{K}N}}
\newcommand{\Kmpp}{\ensuremath{K^{-}pp}}
\newcommand{\KNN}{\ensuremath{\bar{K}NN}}
\newcommand{\Lam}{\ensuremath{\Lambda}}
\newcommand{\Sig}{\ensuremath{\Sigma}}
\newcommand{\Xibar}{\ensuremath{\Xi}}
\newcommand{\Lc}{\ensuremath{\Lambda_c^{+}}}
\newcommand{\gevu}{\,GeV/u}
\newcommand{\mevc}{\,MeV/$c$}

% ---------- Emphasis helpers ----------
\newcommand{\lead}[1]{\textcolor{impblue}{\textbf{#1}}}   % blue bold lead-in
\newcommand{\keypt}[1]{\textbf{#1}}                         % key statement
\newcommand{\keyu}[1]{\uline{\textbf{#1}}}                  % strongest emphasis (use sparingly)
\newcommand{\role}[1]{\textcolor{impsky}{\textbf{[#1]}}}    % authorship/role tag
\newcommand{\tech}[1]{\textmd{#1}}                          % technical term/acronym: keep normal weight even inside bold

% ---------- Blue-bar callout block (for "my contribution" etc.) ----------
\newenvironment{callout}
  {\par\smallskip\noindent\begin{list}{}{%
     \setlength{\leftmargin}{1.1em}\setlength{\rightmargin}{0pt}%
     \setlength{\topsep}{2pt}\setlength{\itemsep}{0pt}}%
     \item[]\color{impink}%
     \def\FrameRule{0pt}}
  {\end{list}\par\smallskip}

% Title banner used at the top of each document
\newcommand{\hjbanner}[3]{% #1 = CN title, #2 = EN subtitle, #3 = meta line
  {\noindent\Huge\bfseries\color{impblue} #1\par}
  \vspace{2pt}
  {\noindent\large\color{impink} #2\par}
  \vspace{4pt}
  {\noindent\small\color{impgray} #3\par}
  \vspace{2pt}
  {\noindent\color{impblue!45}\rule{\linewidth}{1.2pt}\par}
  \vspace{4pt}
}

\usepackage[hidelinks]{hyperref}

%  Compile with: xelatex
% ============================================================

\usepackage{geometry}
\geometry{a4paper, top=2.4cm, bottom=2.3cm, left=2.4cm, right=2.4cm, headsep=10pt}

\usepackage{amsmath,amssymb}
\usepackage{bm}
\usepackage{graphicx}
\usepackage{wrapfig}
\usepackage{float}
\usepackage{array}
\usepackage{booktabs}
\usepackage{tabularx}
\usepackage{enumitem}
\usepackage[dvipsnames,table]{xcolor}
\usepackage{titlesec}
\usepackage{fancyhdr}
\usepackage[normalem]{ulem}
\usepackage{caption}
\usepackage{tikz}
\usetikzlibrary{arrows.meta,positioning,calc,shapes.geometric,backgrounds,fit}
\usepackage{fontspec}

% ---------- Fonts: portable across machines (see README.md) ----------
% Latin = XCharter (bundled with TeX Live). CJK = macOS system fonts when present
% (preserves the original look), else Fandol (bundled with TeX Live) so the document
% builds on ANY TeX Live install (Mac, Linux, Windows) with no extra font setup.
\setmainfont{XCharter}% Latin: XCharter ships with TeX Live (always available)
\IfFontExistsTF{Songti SC}
  {\setCJKmainfont{Songti SC}\newCJKfontfamily\song{Songti SC}}
  {\setCJKmainfont{FandolSong-Regular.otf}[BoldFont=FandolSong-Bold.otf]%
   \newCJKfontfamily\song{FandolSong-Regular.otf}}
\IfFontExistsTF{Heiti SC}
  {\setCJKsansfont{Heiti SC}\newCJKfontfamily\hei{Heiti SC}}% CJK sans for emphasis/headings
  {\setCJKsansfont{FandolHei-Regular.otf}[BoldFont=FandolHei-Bold.otf]%
   \newCJKfontfamily\hei{FandolHei-Regular.otf}}
% route enclosed numbers ①-⑳ and Roman numerals to the CJK font (XCharter lacks them)
\xeCJKDeclareCharClass{CJK}{"2160 -> "2188}
\xeCJKDeclareCharClass{CJK}{"2460 -> "24FF}

% ---------- Colour palette ----------
\definecolor{impblue}{RGB}{31,73,125}     % deep institutional blue
\definecolor{impsky}{RGB}{46,116,181}     % brighter accent blue
\definecolor{impgray}{RGB}{120,120,120}   % header/footer grey
\definecolor{impink}{RGB}{20,40,70}        % near-black ink for key text
\definecolor{impbg}{RGB}{238,243,249}      % pale blue panel background

% ---------- Paragraph style: clean block layout ----------
\setlength{\parindent}{0pt}
\setlength{\parskip}{0.55em}
\linespread{1.06}
\emergencystretch=1.6em   % absorb long unbreakable Latin/math runs, avoid overfull
\setlist{nosep, leftmargin=1.5em, topsep=2pt, itemsep=3pt}

% ---------- Section headings (blue, ruled) ----------
\titleformat{\section}[hang]
  {\Large\bfseries\color{impblue}}{\thesection.}{0.5em}{}
  [{\vspace{1pt}\color{impblue!35}\titlerule[1pt]}]
\titlespacing*{\section}{0pt}{1.1em}{0.5em}

\titleformat{\subsection}[hang]
  {\large\bfseries\color{impsky}}{\thesubsection}{0.5em}{}
\titlespacing*{\subsection}{0pt}{0.8em}{0.25em}

\titleformat{\subsubsection}[runin]
  {\bfseries\color{impink}}{}{0em}{}[\,]
\titlespacing*{\subsubsection}{0pt}{0.4em}{0.5em}

% ---------- Captions ----------
\captionsetup{font=small, labelfont={bf,color=impblue}, labelsep=period, skip=4pt}

% ---------- Running header / footer ----------
\newcommand{\doctitle}{}
\pagestyle{fancy}
\fancyhf{}
\setlength{\headheight}{14pt}
\fancyhead[L]{\small\color{impgray}\,马越\ \textperiodcentered\ Yue Ma}
\fancyhead[R]{\small\color{impgray}\doctitle\,}
\fancyfoot[C]{\small\color{impgray}\thepage}
\renewcommand{\headrulewidth}{0.4pt}
\renewcommand{\footrulewidth}{0pt}
\renewcommand{\headrule}{\hbox to\headwidth{\color{impblue!30}\leaders\hrule height \headrulewidth\hfill}}

% ---------- Physics macros ----------
\newcommand{\hn}[2]{\ensuremath{^{#1}_{\Lambda}\mathrm{#2}}}        % single-Lambda hypernucleus  \hn{3}{H}
\newcommand{\hnn}[2]{\ensuremath{^{#1}_{\Lambda\Lambda}\mathrm{#2}}}% double-Lambda hypernucleus  \hnn{6}{He}
\newcommand{\hsig}[2]{\ensuremath{^{#1}_{\Sigma}\mathrm{#2}}}       % Sigma hypernucleus  \hsig{3}{H}
\newcommand{\dbar}{\ensuremath{\bar{d}}}
\newcommand{\pbar}{\ensuremath{\bar{p}}}
\newcommand{\Kbar}{\ensuremath{\bar{K}}}
\newcommand{\KN}{\ensuremath{\bar{K}N}}
\newcommand{\Kmpp}{\ensuremath{K^{-}pp}}
\newcommand{\KNN}{\ensuremath{\bar{K}NN}}
\newcommand{\Lam}{\ensuremath{\Lambda}}
\newcommand{\Sig}{\ensuremath{\Sigma}}
\newcommand{\Xibar}{\ensuremath{\Xi}}
\newcommand{\Lc}{\ensuremath{\Lambda_c^{+}}}
\newcommand{\gevu}{\,GeV/u}
\newcommand{\mevc}{\,MeV/$c$}

% ---------- Emphasis helpers ----------
\newcommand{\lead}[1]{\textcolor{impblue}{\textbf{#1}}}   % blue bold lead-in
\newcommand{\keypt}[1]{\textbf{#1}}                         % key statement
\newcommand{\keyu}[1]{\uline{\textbf{#1}}}                  % strongest emphasis (use sparingly)
\newcommand{\role}[1]{\textcolor{impsky}{\textbf{[#1]}}}    % authorship/role tag
\newcommand{\tech}[1]{\textmd{#1}}                          % technical term/acronym: keep normal weight even inside bold

% ---------- Blue-bar callout block (for "my contribution" etc.) ----------
\newenvironment{callout}
  {\par\smallskip\noindent\begin{list}{}{%
     \setlength{\leftmargin}{1.1em}\setlength{\rightmargin}{0pt}%
     \setlength{\topsep}{2pt}\setlength{\itemsep}{0pt}}%
     \item[]\color{impink}%
     \def\FrameRule{0pt}}
  {\end{list}\par\smallskip}

% Title banner used at the top of each document
\newcommand{\hjbanner}[3]{% #1 = CN title, #2 = EN subtitle, #3 = meta line
  {\noindent\Huge\bfseries\color{impblue} #1\par}
  \vspace{2pt}
  {\noindent\large\color{impink} #2\par}
  \vspace{4pt}
  {\noindent\small\color{impgray} #3\par}
  \vspace{2pt}
  {\noindent\color{impblue!45}\rule{\linewidth}{1.2pt}\par}
  \vspace{4pt}
}

\usepackage[hidelinks]{hyperref}

% --- Japanese fonts (override hjstyle's Chinese CJK fonts) ---
\setCJKmainfont{Hiragino Mincho ProN}[AutoFakeBold=3]
\setCJKsansfont{Hiragino Sans}[AutoFakeBold=3]
\fancyhead[L]{\small\color{impgray}\,馬越\ \textperiodcentered\ Yue Ma}
\renewcommand{\doctitle}{履歴書}

% date | content row
\newcommand{\cvrow}[2]{%
  \noindent\makebox[2.7cm][l]{\textcolor{impblue}{\small\bfseries #1}}%
  \parbox[t]{\dimexpr\linewidth-2.8cm\relax}{#2}\par\vspace{3pt}}

\begin{document}

\hjbanner{履歴書 \quad\normalsize Curriculum Vitae}
  {Yue Ma（馬越）—— 実験核物理学者}
  {専任研究員（理化学研究所 少数多体系物理研究室）\quad\textbar\quad 生年月日 1979-07-19 \quad\textbar\quad \href{mailto:y.ma@riken.jp}{y.ma@riken.jp} \quad\textbar\quad \href{https://ymaphys.github.io}{ymaphys.github.io} \quad\textbar\quad \href{https://orcid.org/0000-0001-8412-8308}{ORCID 0000-0001-8412-8308}}

% ----------------------------------------------------------------------
\section{概要}

実験核物理学者、理化学研究所 専任研究員（任期なし）。日本（J-PARC、SPring-8、KEK）およびドイツ（GSI、DESY）の大型実験施設において、15 年以上にわたる第一線の研究・指導経験を有し、若手研究者の育成と国際共同研究ネットワークの構築に長く携わってきた。\keypt{現在、\tech{J-PARC E73}（ハイパートリトン寿命）の実験スポークスパーソン、ならびに \tech{SPring-8 LEPS2}（核子短距離相関）および \tech{J-PARC} 反重陽子実験の提案スポークスパーソンを務める。}

% ----------------------------------------------------------------------
\section{学術的指導・役職}
\begin{itemize}
  \item \keypt{\tech{J-PARC E73}、ハイパートリトン寿命 —— 実験スポークスパーソン。} ビームタイムを完了し、\hn{4}{H} の寿命と \hn{3,4}{H} の生成断面積を発表済み。\hn{3}{H} 寿命の解析は進行中。
  \item \keypt{\tech{SPring-8 LEPS2}、核子短距離相関 —— 提案スポークスパーソン。} $\Lambda$ ハイペロンをプローブとして原子核内の核子間短距離相関を研究（進行中）。
  \item \keypt{\tech{J-PARC} 反重陽子実験 —— 提案スポークスパーソン。} 反重陽子–原子核相互作用と対消滅機構（進行中）。
\end{itemize}

\section{主要な研究成果}
\begin{enumerate}
  \item \keypt{\Kmpp（\KNN）反 $K$ 中間子–原子核束縛状態の発見} —— 創設メンバーとしてマルチヒット TDC を開発し、偏極解析により信号を抽出、深く束縛した \Kmpp 状態の世界初の証拠を与えた（J-PARC E15）。\,{\small Phys.\ Lett.\ B \textbf{789} (2019); Phys.\ Rev.\ C \textbf{102} (2020).}
  \item \keypt{\hn{4}{H} 超核寿命の精密測定} —— スポークスパーソン兼責任著者として、高速 $\pi^0$ を単一の高エネルギー $\gamma$ で標識する鍵となるアイデアを提案し測定を主導、$\tau=206\pm8\pm12$\,ps（J-PARC E73）。\,{\small Phys.\ Lett.\ B \textbf{845} (2023); \textbf{873} (2026).}
  \item \keypt{\KN($I{=}1$) 閾値近傍の $\Lambda\pi$ 構造の世界初観測} —— 筆頭著者として Belle データ解析を主導。\,{\small Phys.\ Rev.\ Lett.\ \textbf{130} (2023).}
  \item \keypt{\tech{TES} マイクロカロリメータのハドロン原子分光への世界初適用} —— SDD データ解析とエネルギー較正をサブ eV 精度で実現（J-PARC E62）。\,{\small Prog.\ Theor.\ Exp.\ Phys.\ \textbf{2016}, 091D01; Phys.\ Rev.\ Lett.\ \textbf{128} (2022).}
  \item \keypt{エネルギー分解型ミュオンスピン分光（$\mu$SR）の世界初実証} —— 筆頭著者兼責任著者として、鉛ガラス・チェレンコフカロリメータによるイベント毎の陽電子エネルギー選別を導入して実証実験を主導、ビームタイムを増やすことなく性能指数を 26\% 向上させた（J-PARC MLF）。\,{\small Nucl.\ Instrum.\ Meth.\ A, in press (2026).}
\end{enumerate}

% ----------------------------------------------------------------------
\begin{minipage}[t]{0.49\linewidth}
\section{学歴}
\cvrow{2006--2009}{\textbf{博士（理学）}、東北大学（仙台）}
\cvrow{2004--2006}{修士（理学）、東北大学（仙台）}
\cvrow{1998--2002}{\textbf{学士（核物理学）}、蘭州大学}
\end{minipage}\hfill
\begin{minipage}[t]{0.49\linewidth}
\section{職歴}
\cvrow{2012--現在}{\textbf{専任研究員（任期なし）}、理化学研究所}
\cvrow{2009--2012}{博士研究員、マインツ・ヘルムホルツ研究所（独）}
\cvrow{2002--2004}{研究実習生、中国科学院 近代物理研究所}
\end{minipage}

% ----------------------------------------------------------------------
\begin{minipage}[t]{0.54\linewidth}
\section{競争的研究資金}
\cvrow{2026--2027}{学術変革領域研究（A）公募研究：J-PARC における $\Sigma$ ハイパートリトン（\hsig{3}{H}）束縛状態の探索（肥山詠美子 領域内）。\textbf{研究代表者}}
\cvrow{2024--2028}{基盤研究（A）：光子ビームによる原子核中の $\rho$ 中間子質量変化の測定。研究分担者}
\cvrow{2023--2025}{理化学研究所 奨励課題：物質–反物質相互作用の探索。\textbf{研究代表者}}
\cvrow{2017--2020}{若手研究（A）17H04842：ハイパートリトン寿命パズルの解決。\textbf{研究代表者}}
\end{minipage}\hfill
\begin{minipage}[t]{0.42\linewidth}
\section{学生指導}
\cvrow{2017--2020}{C.\ Han（博士）、中国科学院大学}
\cvrow{2012--2014}{Q.\ Zhang（博士）、蘭州大学}
\cvrow{2010--2011}{D.\ Lin（博士）；Christina、Pascal（修士）、マインツ大学}
\end{minipage}

% ----------------------------------------------------------------------
\begin{minipage}[t]{0.54\linewidth}
\section{学会運営・招待講演}
\cvrow{2019--2021}{理化学研究所 研究者会議運営委員；奨励課題 物理分野コーディネーター／総括コーディネーター}
\cvrow{2022}{招待講演、第 14 回 ハイパー核・ストレンジ粒子物理 国際会議（HYP 2022、プラハ）}
\cvrow{2025}{招待講演、第 15 回 HYP（HYP 2025、東京）}
\end{minipage}\hfill
\begin{minipage}[t]{0.42\linewidth}
\section{教育・語学}
\cvrow{2018/23/24}{招待講師、ハドロン物理学 夏期講習、中国科学院大学}
\cvrow{語学}{中国語（母語）$\cdot$ 英語（流暢）$\cdot$ 日本語（業務遂行レベル）}
\end{minipage}

% ----------------------------------------------------------------------
\section{主要論文 \normalfont\small（主導的役割と影響度の順；著者の役割を併記）}

{\small 発表論文 \keypt{100 編以上}（Belle / Belle\,II 共同研究論文を含む。末尾の節を参照）。以下は筆頭著者・責任著者・スポークスパーソン、または主導的に貢献した代表的成果。}

{\small\color{impgray} 注：ハドロン物理学の論文は姓のアルファベット順に記載される慣例があり、下記の\textbf{筆頭著者}／\textbf{責任著者}／\textbf{スポークスパーソン}の標記は実際の主導的役割を示す。}

{\small
\begin{enumerate}\setlength{\itemsep}{2pt}
  \item \textbf{Y.~Ma} \textit{et al.}, \textit{First observation of $\Lambda\pi^+$ and $\Lambda\pi^-$ signals near the $\bar{K}N(I{=}1)$ mass threshold in $\Lambda_c^+\!\to\Lambda\pi^+\pi^+\pi^-$ decay}, Phys.\ Rev.\ Lett.\ \textbf{130}, 151903 (2023) [Belle]. \role{筆頭著者}
  \item T.~Akaishi \textit{et al.}, \textit{Precise lifetime measurement of $^{4}_{\Lambda}$H hypernucleus using a novel production method}, Phys.\ Lett.\ B \textbf{845}, 138128 (2023). \role{スポークスパーソン \& 責任著者}
  \item T.~Akaishi \textit{et al.}, \textit{Measurement of $^{3,4}$He$(K^-,\pi^0)$ $^{3,4}_{\Lambda}$H reaction cross section and evaluation of hypertriton $\Lambda$ binding energy}, Phys.\ Lett.\ B \textbf{873}, 140163 (2026). \role{スポークスパーソン}
  \item \textbf{Y.~Ma} \textit{et al.}, \textit{Demonstration of energy-resolved muon spin spectroscopy using a Cherenkov calorimeter}, Nucl.\ Instrum.\ Meth.\ A, in press (2026). \role{筆頭著者 \& 責任著者}
  \item \textbf{Y.~Ma} \textit{et al.}, \textit{$\gamma$-ray spectroscopy study of $^{11}_{\Lambda}$B and $^{12}_{\Lambda}$C}, Eur.\ Phys.\ J.\ A \textbf{33}, 243 (2007). \role{筆頭著者}
  \item S.~Ajimura \textit{et al.}, \textit{$K^-pp$, a $\bar{K}$-meson nuclear bound state, observed in $^3$He$(K^-,\Lambda p)n$ reactions}, Phys.\ Lett.\ B \textbf{789}, 620 (2019). \role{主要著者の一人}
  \item T.~Yamaga \textit{et al.}, \textit{Observation of a $\bar{K}NN$ bound state in the $^3$He$(K^-,\Lambda p)n$ reaction}, Phys.\ Rev.\ C \textbf{102}, 044002 (2020). \role{主要貢献者の一人}
  \item T.~Hashimoto \textit{et al.}, \textit{Strong-interaction effects in kaonic-helium isotopes at sub-eV precision with X-ray microcalorimeters}, Phys.\ Rev.\ Lett.\ \textbf{128}, 112503 (2022). \role{主要貢献者の一人}
  \item B.~S.~Henderson \textit{et al.}, \textit{Hard two-photon contribution to elastic lepton-proton scattering (OLYMPUS)}, Phys.\ Rev.\ Lett.\ \textbf{118}, 092501 (2017). \role{主要貢献者の一人；輝度モニターの開発を主導}
  \item K.~Hosomi \textit{et al.}, \textit{Precise determination of $^{12}_{\Lambda}$C level structure by $\gamma$-ray spectroscopy}, Prog.\ Theor.\ Exp.\ Phys.\ \textbf{2015}, 081D01 (2015). \role{第二著者}
\end{enumerate}
}

% ----------------------------------------------------------------------
\section{Belle / Belle\,II 共同研究論文}

{\small Belle / Belle\,II 共同研究グループの正式メンバー（\keypt{2023 年〜}）として、共同研究論文 \keypt{70 編以上}に署名。主題別の代表的論文：}

{\small\setlength{\emergencystretch}{3em}

\smallskip
\noindent\lead{I.\ $B$ 中間子の CP 対称性の破れと CKM 行列要素}\quad{\color{impgray}時間依存 CP 対称性の破れ（$\sin 2\phi_1$ など）、$|V_{cb}|$／$|V_{ub}|$、CKM 角 $\phi_3$。}
\begin{enumerate}\setlength{\itemsep}{1.5pt}
  \item \textit{Measurement of $\sin 2\phi_{1}$ in $B^{0}\!\to J/\psi K_{S}^{0}$ decays with a graph-neural-network flavor tagger}, Phys.\ Rev.\ D \textbf{110}, 012001 (2024)\;[Belle\,II].
  \item \textit{Measurement of CP violation in $B^{0}\!\to K_{S}^{0}\pi^{0}$ decays}, Phys.\ Rev.\ Lett.\ \textbf{131}, 111803 (2023)\;[Belle\,II].
  \item \textit{Determination of $|V_{cb}|$ using $\bar{B}^{0}\!\to D^{*+}\ell^{-}\bar\nu_{\ell}$ decays}, Phys.\ Rev.\ D \textbf{108}, 092013 (2023)\;[Belle\,II].
  \item \textit{Measurement of inclusive $B\!\to X_{u}\ell\nu$ partial branching fractions and $|V_{ub}|$}, Phys.\ Rev.\ D \textbf{113}, 032004 (2026)\;[Belle\,II].
  \item \textit{Determination of the CKM angle $\phi_{3}$ from a combination of Belle and Belle\,II results}, J.\ High Energy Phys.\ \textbf{10} (2024) 143 [Belle/Belle\,II].
\end{enumerate}

\smallskip
\noindent\lead{II.\ 稀崩壊とレプトン普遍性（新物理探索）}\quad{\color{impgray}$B^+\!\to K^+\nu\bar\nu$ の世界初の証拠、$R(D^{*})$ などのレプトン普遍性検証、$\tau$ レプトンフレーバー非保存の探索。}
\begin{enumerate}\setlength{\itemsep}{1.5pt}
  \item \textit{Evidence for $B^{+}\!\to K^{+}\nu\bar\nu$ decays}, Phys.\ Rev.\ D \textbf{109}, 112006 (2024)\;[Belle\,II].
  \item \textit{Test of light-lepton universality with a measurement of $R(D^{*})$ using hadronic $B$ tagging}, Phys.\ Rev.\ D \textbf{110}, 072020 (2024)\;[Belle\,II].
  \item \textit{First measurement of $R(X_{\tau/\ell})$ as an inclusive test of the $b\!\to c\tau\nu$ anomaly}, Phys.\ Rev.\ Lett.\ \textbf{132}, 211804 (2024)\;[Belle\,II].
  \item \textit{Tests of light-lepton universality in angular asymmetries of $B^{0}\!\to D^{*-}\ell\nu$ decays}, Phys.\ Rev.\ Lett.\ \textbf{131}, 181801 (2023)\;[Belle\,II].
  \item \textit{Measurement of the $B^{+}\!\to\tau^{+}\nu_{\tau}$ branching fraction with a hadronic tagging method}, Phys.\ Rev.\ D \textbf{112}, 072002 (2025)\;[Belle\,II].
  \item \textit{First search for $B\!\to X_{s}\nu\bar\nu$ decays}, Phys.\ Rev.\ Lett.\ \textbf{136}, 231801 (2026)\;[Belle\,II].
  \item \textit{Search for lepton-flavor-violating $\tau$ decays to $\ell\alpha$}, J.\ High Energy Phys.\ \textbf{08} (2025) 155 [Belle].
  \item \textit{Search for a $\tau^{+}\tau^{-}$ resonance in $e^{+}e^{-}\!\to\mu^{+}\mu^{-}\tau^{+}\tau^{-}$ events}, Phys.\ Rev.\ Lett.\ \textbf{131}, 121802 (2023)\;[Belle\,II].
\end{enumerate}

\smallskip
\noindent\lead{III.\ チャーム／ストレンジ・バリオン分光とエキゾチックハドロン}{\color{impgray}（自身の研究と密接に関連）}\quad{\color{impgray}$\Xi_c$ の崩壊と CP、$D_{s0}^{*}(2317)$ 輻射崩壊の世界初観測、$P_{cs}$ ペンタクォーク、$\Xi^0 p$／$\Omega^- p$／$\Omega^- n$ 二重バリオンの探索。}
\begin{enumerate}\setlength{\itemsep}{1.5pt}
  \item \textit{Observation of the radiative decay $D_{s0}^{*}(2317)^{+}\!\to D_{s}^{*+}\gamma$}, Phys.\ Rev.\ Lett.\ \textbf{136}, 241901 (2026) [Belle/Belle\,II].
  \item \textit{Search for $P_{cs}(4459)$ and $P_{cs}(4338)$ in $\Upsilon(1S,2S)$ inclusive decays}, Phys.\ Rev.\ Lett.\ \textbf{135}, 041901 (2025) [Belle/Belle\,II].
  \item \textit{Search for CP violation in $\Xi_{c}^{+}\!\to\Sigma^{+}h^{+}h^{-}$ and $\Lambda_{c}^{+}\!\to p\,h^{+}h^{-}$ decays}, Phys.\ Rev.\ D \textbf{113}, 032017 (2026)\;[Belle\,II].
  \item \textit{Observation of the singly Cabibbo-suppressed decays $\Xi_{c}^{+}\!\to pK_{S}^{0}$, $\Lambda\pi^{+}$, $\Sigma^{0}\pi^{+}$}, J.\ High Energy Phys.\ \textbf{03} (2025) 061 [Belle/Belle\,II].
  \item \textit{Measurement of branching fractions of $\Xi_{c}^{0}\!\to\Xi^{0}\pi^{0},\,\Xi^{0}\eta,\,\Xi^{0}\eta'$}, J.\ High Energy Phys.\ \textbf{10} (2024) 045 [Belle/Belle\,II].
  \item \textit{Observation of $B^{+}\!\to\Sigma_{c}(2455)^{++}\bar\Xi_{c}^{-}$ and $B^{0}\!\to\Sigma_{c}(2455)^{0}\bar\Xi_{c}^{0}$}, Phys.\ Rev.\ D \textbf{112}, L051101 (2025)\;[Belle/Belle\,II].
\end{enumerate}
}
\end{document}
